% Ad Familiares 2.16 (ad-familiares-02-16) --- English translation
% Translated by Claude (Anthropic) from the Latin (Perseus).
% Style guide: STYLE.md.

\ciceroLetterOpener{M. CICERO IMP. S. D. M. CAELIO.}

\ciceroSection{1} Your letter would have caused me great pain, were it not that reason itself by now has driven out every distress, and that long despair over public affairs has hardened my mind against any fresh grief. Yet how it came about that you should have inferred from my earlier letter what you write you did, I do not know. What was there in it but a lament for the times? --- a lament which troubles my mind no more than it does yours. For I do not recognize the keenness of intellect which I myself observe in you as failing to see what I see; what amazes me is that you, who ought to know me through and through, could have been brought to suppose me either so improvident as to break away from the rising fortune for the failing and the all-but-prostrate, or so inconstant as to fling away the goodwill I had won from a man at the height of his power, to desert myself, and to take part --- a thing I avoided at the outset and have always avoided --- in a civil war. What then is this ``grim decision'' of mine?

\ciceroSection{2} To withdraw, perhaps, into some quiet retreat. For you know my temper --- of which you once had something like its match --- and you know besides my disgust, in the eyes as much as in the stomach, at the insolence of the overbearing. To this is added the awkward pomp of my lictors and the name of \textit{imperator} by which I am still addressed. If I were rid of that load, I should be content with however small a hiding-place in Italy; but this laurel of mine catches the eye, and now the busy little tongues, of the malevolent. And though that was the case, still I have never given a thought to leaving the country except with your good people's approval. But you know my modest estates: I am obliged to be on one of them, lest I become a burden to my friends. And because I am at my easiest on the coastal ones, I raise in some quarters the suspicion that I mean to take ship --- which I might not, in fact, be unwilling to do, if it could be done for the sake of quiet; but for war, how does that fit? --- particularly to make war on a man whom I hope I have satisfied, on the side of a man who in no way can ever now be satisfied.

\ciceroSection{3} Then again, you could most easily have made out my view back at the time when you met me on the road at Cumae. For I did not hide T.\ Ampius's report from you; you saw how I shrank from the abandonment of the city when I heard of it. Did I not declare to you that I would endure anything sooner than leave Italy for civil war? What then has happened to make me change my mind? Has not everything rather happened to keep me in it? Believe this of me, as I think you do: out of all these miseries I seek nothing else than that men should one day understand that I have wanted nothing more than peace, and that, when peace had been despaired of, I have shunned nothing so much as civil arms. I do not think I shall ever regret this constancy. I remember that our friend Q.\ Hortensius used to make it his boast, in this same line, that he had never taken part in a civil war; my own credit will shine the brighter, because in his case the thing was set down to want of spirit, while in mine I do not think such a verdict is open.

\ciceroSection{4} Nor do those prospects frighten me which, in the most loyal and affectionate spirit, you set out to me as cause for alarm. For there is no bitterness which does not seem to be hanging over us all amid this upheaval of the whole world --- a bitterness which, for the state's sake, I should most gladly have bought off, even at the price of just those private and domestic losses against which you bid me take precautions.

\ciceroSection{5} For my son --- I am glad you hold him dear --- I shall leave, if there is any state to leave it in, patrimony enough in the memory of my name; if there is none, then nothing will befall him separately from the rest of the citizens. As for your urging me to look out for my son-in-law, that excellent young man and one most dear to me --- can you doubt, when you know how highly I value him, and how much higher still my Tullia, that this care presses on me with the greatest weight? all the more because amid our common miseries this was the one gleam of hope I cherished, that my Dolabella, or rather ours, would be set free from the troubles he had brought on himself by his generosity. I should be glad if you would inquire what days he endured while he was in the city --- how bitter they were to him, how dishonourable to me, his father-in-law.

\ciceroSection{6} And so I am neither awaiting this Spanish outcome --- of which, by the way, I am as well informed as you write you are --- nor laying any subtle scheme. If there is ever to be a commonwealth again, surely there will be a place for me in it; if not, you yourself, I imagine, will come into the same wildernesses where you will hear that I have settled. Yet perhaps I am playing the prophet, and all this will end better. For I remember the despair of the men who were old when I was young. Maybe I am imitating them now and falling into the failing of age. I should be glad if it were so; but still ---.

\ciceroSection{7} I expect you have heard that I have had the \textit{toga praetexta} woven for Oppius's boy; our friend Curtius is thinking of dipping his twice in dye, but the dyer holds him up. I have thrown that in by way of showing you that I can still find something to laugh at in the midst of vexation. About Dolabella's affair, the line I have written --- look to it, please, as though it were your own. The last word will be this: I shall do nothing turbulent, nothing rash; but I do beg you, wherever I may be on earth, to look after me and my children as our friendship and your own loyalty demand.
